% Part: first-order-logic 
% Chapter: axiomatic-deduction 
% Section: deduction-theorem

% verification of properties of provability needed for maximally 
% consistent sets in the completeness chapter.

\documentclass[../../../include/open-logic-section]{subfiles}

\begin{document}

\iftag{FOL}
      {\olfileid[pt]{fol}{axd}{ded}}
      {\olfileid[pt]{pl}{axd}{ded}}

\olsection{O teorema da dedução}

Como vimos, fornecer !!{derivation}s em um sistema axiomático é
trabalhoso, e pode ser difícil encontrar !!{derivation}s. Em vez de
escrever efetivamente longas listas de !!{formula}s, em geral é mais
fácil argumentar que tais !!{derivation}s existem usando alguns
resultados simples. Já estabelecemos três resultados desse tipo:
\olref[ptn]{prop:reflexivity} diz que sempre podemos afirmar $\Gamma
\Proves !A$ quando sabemos que $!A \in \Gamma$.
\olref[ptn]{prop:monotonicity} diz que, se $\Gamma \Proves !A$, então
também $\Gamma \cup \{!B\} \Proves !A$. E
\olref[ptn]{prop:transitivity} implica que, se $\Gamma \Proves !A$ e
$!A \Proves !B$, então $\Gamma \Proves !B$. Eis outro resultado
simples, uma versão ``meta'' do modus ponens:

\begin{prop}
  \ollabel{prop:mp} Se $\Gamma \Proves !A$ e $\Gamma \Proves !A \lif
  !B$, então $\Gamma \Proves !B$.
\end{prop}

\begin{proof}
  Temos $\{!A, !A \lif !B\} \Proves !B$:
  \begin{derivation}
    1. & $!A$ & Hip.\\
    2. & $!A \lif !B$ & Hip.\\
    3. & $!B$ & 1, 2, MP
  \end{derivation}
  Por \olref[ptn]{prop:transitivity}, $\Gamma \Proves !B$.
\end{proof}

O resultado mais importante que usaremos neste contexto é o teorema da
dedução:

\begin{thm}[Teorema da dedução]
\ollabel{thm:deduction-thm} $\Gamma \cup \{!A\} \Proves !B$ se, e
somente se, $\Gamma \Proves !A \lif !B$.
\end{thm}

\begin{proof}
A direção ``se'' é imediata. Se $\Gamma \Proves !A \lif !B$, então
também $\Gamma \cup \{!A\} \Proves !A \lif !B$ por
\olref[ptn]{prop:monotonicity}. Além disso,
$\Gamma \cup \{!A\} \Proves !A$ por \olref[ptn]{prop:reflexivity}.
Logo, por \olref{prop:mp}, $\Gamma \cup \{!A\} \Proves !B$.

Para a direção ``somente se'', procedemos por indução sobre o
comprimento da !!{derivation} de $!B$ a partir de
$\Gamma \cup \{!A\}$.

Na base da indução, demonstramos a afirmação para toda !!{derivation} de
comprimento~$1$. !!^a{derivation} de~$!B$ a partir de
$\Gamma \cup \{!A\}$ com comprimento~$1$ consiste apenas em $!B$; se
ela está correta, $!B$ ou está $\in \Gamma \cup \{!A\}$ ou é um axioma.
Se $!B \in \Gamma$ ou é um axioma, então $\Gamma \Proves !B$. Também
temos $\Gamma \Proves !B \lif (!A \lif !B)$ por
\olref[prp]{ax:lif1}, e \olref{prop:mp} fornece
$\Gamma \Proves !A \lif !B$. Se $!B \in \{ !A\}$, então
$\Gamma \Proves !A \lif !B$, pois a última !!{sentence}~$!A \lif !B$ é
a mesma que $!A \lif !A$, e já !!{derive}d isso em
\olref[pro]{ex:identity}.

No passo de indução, suponha que !!a{derivation} de~$!B$ a partir de
$\Gamma \cup \{!A\}$ termine em um passo~$!B$ justificado por modus
ponens. (Se ele não for justificado por modus ponens, $!B \in \Gamma$,
$!B \ident !A$ ou $!B$ será um axioma, e se aplicará o mesmo raciocínio
da base da indução.) Então alguns passos anteriores da !!{derivation}
são $!C \lif !B$ e $!C$, para alguma !!{formula}~$!C$, isto é,
$\Gamma \cup \{!A\} \Proves !C \lif !B$ e $\Gamma \cup \{!A\}
\Proves !C$, e as respectivas !!{derivation}s são mais curtas; portanto,
a hipótese de indução se aplica a elas. Assim, temos ambos:
\begin{align*}
  & \Gamma \Proves !A \lif (!C \lif !B); \\
  & \Gamma \Proves !A \lif !C.
\end{align*}
Mas também
\[
\Gamma \Proves (!A \lif (!C \lif !B)) \lif
((!A\lif !C)  \lif (!A \lif !B)),
\]
por \olref[prp]{ax:lif2}, e duas aplicações de \olref{prop:mp} fornecem
$\Gamma \Proves !A \lif !B$, como desejado.
\end{proof}

Observe como \olref[prp]{ax:lif1} e \olref[prp]{ax:lif2} foram escolhidos
precisamente para que o teorema da dedução fosse válido.

Apresentamos a seguir alguns fatos úteis sobre !!{derivability}, que
deixamos como exercícios.

\begin{prop}
  \ollabel{prop:derivfacts}
\begin{enumerate}
\item $\Proves (!A \lif !B) \lif ((!B \lif !C)
  \lif (!A \lif !C)$; \ollabel{derivfacts:a}
\item Se $\Gamma \cup \{ \lnot !A\}
  \Proves \lnot !B$, então $\Gamma \cup \{ !B\} \Proves
  !A$ (Contraposição); \ollabel{derivfacts:b}
\item  $\{ !A, \lnot!A\} \Proves
    !B$ (Ex Falso Quodlibet, explosão); \ollabel{derivfacts:c}
\item  $\{ \lnot\lnot!A\} \Proves
  !A$ (Eliminação da dupla negação);\ollabel{derivfacts:d}
\item Se $\Gamma \Proves \lnot\lnot!A$, então $\Gamma \Proves
  !A$;\ollabel{derivfacts:e}
\end{enumerate}
\end{prop}

\tagprob{FOL}
\begin{prob}
  Demonstre \olref[fol][axd][ded]{prop:derivfacts}
\end{prob}
\tagendprob

\tagprob{notFOL}
\begin{prob}
  Demonstre \olref[pl][axd][ded]{prop:derivfacts}
\end{prob}
\tagendprob

\end{document}
